\documentclass[11pt]{article}
\usepackage[margin=1in]{geometry}
\usepackage{xcolor}
\usepackage{tcolorbox}
\usepackage{hyperref}
\usepackage{enumitem}
\usepackage{amsmath}
\usepackage{booktabs}

% Define colors
\definecolor{osaorange}{RGB}{220,115,38}
\definecolor{correctgreen}{RGB}{34,139,34}

% Configure hyperref
\hypersetup{
    colorlinks=true,
    linkcolor=blue,
    urlcolor=blue
}

\title{}
\author{}
\date{}

\begin{document}

% Header box
\begin{tcolorbox}[colback=osaorange,colframe=black,boxrule=2pt,arc=0pt,left=3pt,right=3pt,top=6pt,bottom=6pt,boxsep=2pt]
\centering
\textcolor{white}{\textbf{\Large H524 Introduction to Biostatistics}}\\
\textcolor{white}{\textbf{\Large Week 5 Assignment - ANSWER KEY}}\\
\textcolor{white}{\textbf{\Large Fall 2025}}
\end{tcolorbox}

\vspace{8pt}

\noindent\textbf{Total Points:} 65 (+ 5 bonus)\\
\textbf{Instructor Use Only}

\section*{Answer Key - Quick Reference}

\subsection*{Part 1: Hypothesis Testing Fundamentals (18 points)}

\begin{table}[h]
\begin{tabular}{clcp{8cm}}
\toprule
\textbf{Q\#} & \textbf{Answer} & \textbf{Pts} & \textbf{Topic} \\
\midrule
1 & \textcolor{correctgreen}{D} & 3 & Null hypothesis definition \\
2 & \textcolor{correctgreen}{C} & 3 & Alternative hypothesis (one-sided, lower) \\
3 & \textcolor{correctgreen}{C} & 4 & P-value interpretation \\
4 & \textcolor{correctgreen}{B} & 2 & Decision rule \\
5 & \textcolor{correctgreen}{A} & 3 & T-statistic = 1.25 \\
6 & \textcolor{correctgreen}{A} & 3 & P-value = 0.2234 \\
\bottomrule
\end{tabular}
\end{table}

\subsection*{Part 2: Type I and Type II Errors (14 points)}

\begin{table}[h]
\begin{tabular}{clcp{8cm}}
\toprule
\textbf{Q\#} & \textbf{Answer} & \textbf{Pts} & \textbf{Topic} \\
\midrule
7 & \textcolor{correctgreen}{C} & 3 & Type I error (false positive) \\
8 & \textcolor{correctgreen}{B} & 3 & Type II error (false negative) \\
9 & \textcolor{correctgreen}{D} & 3 & Error trade-off \\
10 & \textcolor{correctgreen}{B} & 2 & Statistical power \\
11a & \textcolor{correctgreen}{5.18 or 5.181} & 5 & Blood pressure t-statistic \\
11b & \textcolor{correctgreen}{0.0001 or 3.17e-05} & 5 & Blood pressure p-value \\
\bottomrule
\end{tabular}
\end{table}

\subsection*{Part 3: Hypothesis Testing with R (10 points)}

\begin{table}[h]
\begin{tabular}{clcp{8cm}}
\toprule
\textbf{Q\#} & \textbf{Answer} & \textbf{Pts} & \textbf{Topic} \\
\midrule
12a & \textcolor{correctgreen}{-2.71 or -2.7105} & 5 & Cholesterol t-statistic \\
12b & \textcolor{correctgreen}{0.0050 or 0.004967} & 5 & Cholesterol p-value \\
\bottomrule
\end{tabular}
\end{table}

\subsection*{Part 4: Two-Sample Tests (14 points)}

\begin{table}[h]
\begin{tabular}{clcp{8cm}}
\toprule
\textbf{Q\#} & \textbf{Answer} & \textbf{Pts} & \textbf{Topic} \\
\midrule
13 & \textcolor{correctgreen}{B} & 3 & Two-sample application \\
14 & \textcolor{correctgreen}{A} & 4 & Two-sample calculation \\
15 & \textcolor{correctgreen}{D} & 2 & R code for two-sample test \\
\bottomrule
\end{tabular}
\end{table}

\newpage

\section*{Detailed Solutions}

\subsection*{Part 1: Hypothesis Testing Fundamentals}

\textbf{Question 1 (3 points):} Null Hypothesis Definition

\textcolor{correctgreen}{\textbf{Answer: D}}

$H_0$ represents the status quo or no effect; it is the hypothesis we assume true until we have sufficient evidence to reject it. This is the fundamental principle of hypothesis testing - we start by assuming no effect and require strong evidence to reject this assumption.

\vspace{0.5cm}

\textbf{Question 2 (3 points):} Alternative Hypothesis Forms

\textcolor{correctgreen}{\textbf{Answer: C}}

Since the researcher wants to test if the drug LOWERS blood pressure compared to 140 mmHg, the appropriate alternative hypothesis is $H_a$: $\mu < 140$ (one-sided, lower tail). This is a directional hypothesis testing for a decrease.

\vspace{0.5cm}

\textbf{Question 3 (4 points):} P-value Interpretation

\textcolor{correctgreen}{\textbf{Answer: C}}

The p-value is the probability of observing data as extreme as (or more extreme than) what we observed, assuming the null hypothesis is true. It is NOT the probability that $H_0$ is true given our data.

\vspace{0.5cm}

\textbf{Question 4 (2 points):} Decision Rule

\textcolor{correctgreen}{\textbf{Answer: B}}

Since p-value (0.032) $< \alpha$ (0.05), we reject $H_0$. We have sufficient evidence against the null hypothesis at the 5\% significance level.

\vspace{0.5cm}

\textbf{Question 5 (3 points):} T-statistic Calculation

\textcolor{correctgreen}{\textbf{Answer: A (t = 1.25)}}

\textbf{Solution:}
$$t = \frac{\bar{x} - \mu_0}{s/\sqrt{n}} = \frac{52 - 50}{8/\sqrt{25}} = \frac{2}{8/5} = \frac{2}{1.6} = 1.25$$

\vspace{0.5cm}

\textbf{Question 6 (3 points):} P-value from t-statistic

\textcolor{correctgreen}{\textbf{Answer: A (p-value = 0.2234)}}

\textbf{Solution:}
For $t = 1.25$ with $df = 24$ (two-sided test):
\begin{verbatim}
2 * pt(1.25, df = 24, lower.tail = FALSE)
# Result: 0.2234
\end{verbatim}

Since this is a two-sided test, we need to multiply the one-tail probability by 2.

\newpage

\subsection*{Part 2: Type I and Type II Errors}

\textbf{Question 7 (3 points):} Type I Error

\textcolor{correctgreen}{\textbf{Answer: C}}

Type I error occurs when we reject the null hypothesis when it is actually true (false positive). The probability of Type I error is controlled by the significance level $\alpha$.

\vspace{0.5cm}

\textbf{Question 8 (3 points):} Type II Error

\textcolor{correctgreen}{\textbf{Answer: B}}

This is a Type II error ($\beta$) - failing to reject a false null hypothesis. The drug is actually effective (null is false), but the test failed to detect it (failed to reject the null).

\vspace{0.5cm}

\textbf{Question 9 (3 points):} Error Trade-off

\textcolor{correctgreen}{\textbf{Answer: D}}

When you decrease $\alpha$ (become more conservative about rejecting $H_0$), you increase $\beta$ (the probability of failing to reject when you should). Making it harder to reject $H_0$ increases the chance of missing a true effect.

\vspace{0.5cm}

\textbf{Question 10 (2 points):} Statistical Power

\textcolor{correctgreen}{\textbf{Answer: B}}

Power = $1 - \beta$; it's the probability of correctly rejecting a false null hypothesis (detecting a true effect when it exists). Higher power means better ability to detect true effects.

\vspace{0.5cm}

\textbf{Question 11 (10 points):} Blood Pressure Hypothesis Test

\textbf{Data:}
\begin{verbatim}
bp_data <- c(118, 124, 130, 128, 122, 135, 119, 126, 132, 120,
             125, 129, 121, 127, 123, 131, 120, 128, 124, 126)
\end{verbatim}

\textbf{R Code:}
\begin{verbatim}
result <- t.test(bp_data, mu = 120, alternative = "greater")
result$statistic  # t-statistic
result$p.value    # p-value
\end{verbatim}

\textcolor{correctgreen}{\textbf{Answers:}}
\begin{enumerate}[label=\alph*)]
\item \textbf{t-statistic:} \textcolor{correctgreen}{5.18} (or 5.181 for more precision)
\item \textbf{p-value:} \textcolor{correctgreen}{0.0001} (or 3.17e-05 for exact value)
\end{enumerate}

\textbf{Interpretation:} With such a small p-value (much less than 0.05), we have very strong evidence that the mean blood pressure is significantly higher than 120 mmHg.

\newpage

\subsection*{Part 3: Hypothesis Testing with R}

\textbf{Question 12 (10 points):} Cholesterol Hypothesis Test

\textbf{Given:} $n = 40$, $\bar{x} = 185$ mg/dL, $s = 35$ mg/dL

\textbf{Hypotheses:} $H_0$: $\mu = 200$ vs $H_a$: $\mu < 200$ (one-sided, lower tail)

\textbf{Manual Calculation:}
$$t = \frac{\bar{x} - \mu_0}{s/\sqrt{n}} = \frac{185 - 200}{35/\sqrt{40}} = \frac{-15}{35/6.325} = \frac{-15}{5.534} = -2.71$$

\textbf{R Code:}
\begin{verbatim}
# Manual calculation
n <- 40
xbar <- 185
s <- 35
mu0 <- 200

t_stat <- (xbar - mu0) / (s / sqrt(n))
p_value <- pt(t_stat, df = n - 1)

# Or using t.test
# Note: We don't have the actual data, so we use manual calculation
# If we had data: t.test(chol_data, mu = 200, alternative = "less")
\end{verbatim}

\textcolor{correctgreen}{\textbf{Answers:}}
\begin{enumerate}[label=\alph*)]
\item \textbf{t-statistic:} \textcolor{correctgreen}{-2.71} (or -2.7105 for more precision)
\item \textbf{p-value:} \textcolor{correctgreen}{0.0050} (or 0.004967 for more precision)
\end{enumerate}

\textbf{Interpretation:} At $\alpha = 0.01$, the p-value (0.0050) is less than 0.01, so we reject $H_0$. We have strong evidence that the drug reduces cholesterol to below 200 mg/dL on average. The company's claim is supported by the data.

\newpage

\subsection*{Part 4: Two-Sample Tests}

\textbf{Question 13 (3 points):} Two-Sample Test Application

\textcolor{correctgreen}{\textbf{Answer: B}}

Comparing mean recovery times between patients receiving Drug A versus Drug B requires a two-sample t-test because we're comparing two independent groups. The other options all compare a sample mean to a known value (requiring one-sample tests).

\vspace{0.5cm}

\textbf{Question 14 (4 points):} Two-Sample T-test

\textcolor{correctgreen}{\textbf{Answer: A}}

\textbf{Given:}
\begin{itemize}
\item Treatment A: $n_1 = 30$, $\bar{x}_1 = 145$, $s_1 = 15$
\item Treatment B: $n_2 = 35$, $\bar{x}_2 = 138$, $s_2 = 18$
\end{itemize}

\textbf{R Code:}
\begin{verbatim}
# Create data from summary statistics (approximation)
group_a <- rnorm(30, mean = 145, sd = 15)
group_b <- rnorm(35, mean = 138, sd = 18)

# Two-sample t-test
result <- t.test(group_a, group_b, var.equal = FALSE)
result$statistic
result$p.value
\end{verbatim}

\textbf{Result:} $t \approx 1.71$, p-value $\approx 0.0922$

\textbf{Decision:} Since p-value (0.0922) $> \alpha$ (0.05), we fail to reject $H_0$. There is insufficient evidence to conclude that the two treatments have different mean values at the 5\% significance level.

\vspace{0.5cm}

\textbf{Question 15 (2 points):} R Code for Two-Sample Test

\textcolor{correctgreen}{\textbf{Answer: D}}

The correct syntax is \texttt{t.test(group\_a, group\_b)}. This performs a two-sample t-test comparing the means of the two groups.

\newpage

\section*{Bonus Question - Grading Guidelines}

\subsection*{Question 16: AI Verification and Learning (5 bonus points)}

\textcolor{correctgreen}{\textbf{Grading Notes:}}

\textbf{Part a: (1 point)} Student should provide actual AI prompt and response
\begin{itemize}
\item Clear, specific prompt: 0.5 points
\item AI response documented: 0.5 points
\end{itemize}

\textbf{Part b: (2 points)} Student should verify AI answer step-by-step
\begin{itemize}
\item Step-by-step verification with calculations: 1 point
\item Correctly identifying accuracy/errors: 1 point
\end{itemize}

\textbf{Part c: (2 points)} Student should evaluate AI's explanation of one-sided vs. two-sided tests
\begin{itemize}
\item Providing AI's explanation: 0.5 points
\item Thoughtful evaluation of quality and clarity: 1 point
\item Discussion of strengths/weaknesses: 0.5 points
\end{itemize}

\textbf{Look for:}
\begin{itemize}
\item Clear documentation of AI interaction
\item Critical thinking about AI responses
\item Understanding of hypothesis testing concepts
\item Ability to verify and evaluate AI output
\item Recognition that AI can make mistakes
\end{itemize}

\section*{Grading Summary}

\begin{table}[h]
\centering
\begin{tabular}{lcc}
\toprule
\textbf{Section} & \textbf{Points} & \textbf{Score} \\
\midrule
Part 1: Fundamentals (Q1-Q6) & 18 & \\
Part 2: Errors and Power (Q7-Q11) & 14 & \\
Part 3: R Implementation (Q12) & 10 & \\
Part 4: Two-Sample Tests (Q13-Q15) & 14 & \\
\midrule
\textbf{Subtotal} & \textbf{65} & \\
Bonus: AI Learning (Q16) & +5 & \\
\midrule
\textbf{Total Possible} & \textbf{70} & \\
\bottomrule
\end{tabular}
\end{table}

\section*{Common Student Errors}

\begin{itemize}
\item \textbf{Q5:} Forgetting to divide by $\sqrt{n}$ in the denominator
\item \textbf{Q6:} Using one-tail p-value instead of two-tail (forgetting to multiply by 2)
\item \textbf{Q11:} Using wrong alternative (``two.sided'' instead of ``greater'')
\item \textbf{Q12:} Using wrong alternative (``greater'' instead of ``less'')
\item \textbf{Q14:} Attempting to reject at $\alpha = 0.05$ when p-value is close (0.0922)
\end{itemize}

\section*{Note to Graders}

\textbf{Computational Verification:}
\begin{itemize}
\item All numerical answers in this key have been verified using R (see \texttt{verify\_answers.R} in this directory)
\item For Questions 11 and 12: Accept answers within reasonable rounding error
\begin{itemize}
\item Q11a: Accept 5.18, 5.181, or 5.2
\item Q11b: Accept 0.0001, 3.17e-05, or $<$0.001
\item Q12a: Accept -2.71, -2.7105, or -2.7
\item Q12b: Accept 0.005, 0.00497, or 0.0050
\end{itemize}
\end{itemize}

\textbf{Partial Credit Guidelines:}
\begin{itemize}
\item \textbf{Q11 \& Q12 (Computational):} Award partial credit for correct approach even if calculation is wrong
\begin{itemize}
\item Correct code structure with minor errors: 70\% credit
\item Correct formula but computational mistake: 80\% credit
\item Correct method but wrong R function arguments: 60\% credit
\end{itemize}
\item \textbf{Bonus Question (Q16):} Use detailed rubric on pages 5-6; look for critical thinking, not just AI copy-paste
\end{itemize}

\textbf{Academic Integrity:}
\begin{itemize}
\item Watch for identical code/answers across students (especially Q11-12)
\item Bonus question should show student's OWN evaluation, not just AI output
\item Flag suspicious patterns for instructor review
\end{itemize}

\textbf{Common Mistakes to Watch For:}
\begin{itemize}
\item Using \texttt{alternative = "two.sided"} when problem specifies one-sided test
\item Confusing one-tail and two-tail p-values (Q6)
\item Wrong direction for alternative hypothesis (Q11 should be "greater", Q12 should be "less")
\end{itemize}

\end{document}
